\documentclass[10pt,a4paper]{amsart}
\usepackage{amsmath,amssymb,amsthm,graphicx}
\usepackage{verbatim}
\newenvironment{reference}{\comment}{\endcomment}
%\newenvironment{reference}{}{}
\newenvironment{slogan}{\comment}{\endcomment}
\newenvironment{history}{\comment}{\endcomment}

% For commutative diagrams we use Xy-pic
\usepackage[all]{xy}

% We use 2cell for 2-commutative diagrams.
\xyoption{2cell}
% Theorem environments.
%
\theoremstyle{plain}
\newtheorem{theorem}[subsection]{Theorem}
\newtheorem{proposition}[subsection]{Proposition}
\newtheorem{lemma}[subsection]{Lemma}

\theoremstyle{definition}
\newtheorem{definition}[subsection]{Definition}
\newtheorem{example}[subsection]{Example}
\newtheorem{exercise}[subsection]{Exercise}
\newtheorem{situation}[subsection]{Situation}

\theoremstyle{remark}
\newtheorem{remark}[subsection]{Remark}
\newtheorem{remarks}[subsection]{Remarks}

\numberwithin{equation}{subsection}

% Macros
%
\def\lim{\mathop{\mathrm{lim}}\nolimits}
\def\colim{\mathop{\mathrm{colim}}\nolimits}
\def\Spec{\mathop{\mathrm{Spec}}}
\def\Hom{\mathop{\mathrm{Hom}}\nolimits}
\def\Ext{\mathop{\mathrm{Ext}}\nolimits}
\def\SheafHom{\mathop{\mathcal{H}\!\mathit{om}}\nolimits}
\def\SheafExt{\mathop{\mathcal{E}\!\mathit{xt}}\nolimits}
\def\Sch{\mathit{Sch}}
\def\Mor{\mathop{\mathrm{Mor}}\nolimits}
\def\Ob{\mathop{\mathrm{Ob}}\nolimits}
\def\Sh{\mathop{\mathit{Sh}}\nolimits}
\def\NL{\mathop{N\!L}\nolimits}
\def\CH{\mathop{\mathrm{CH}}\nolimits}
\def\proetale{{pro\text{-}\acute{e}tale}}
\def\etale{{\acute{e}tale}}
\def\QCoh{\mathit{QCoh}}
\def\Ker{\mathop{\mathrm{Ker}}}
\def\Im{\mathop{\mathrm{Im}}}
\def\Coker{\mathop{\mathrm{Coker}}}
\def\Coim{\mathop{\mathrm{Coim}}}

% Boxtimes
%
\DeclareMathSymbol{\boxtimes}{\mathbin}{AMSa}{"02}

%
% Macros for moduli stacks/spaces
%
\def\QCohstack{\mathcal{QC}\!\mathit{oh}}
\def\Cohstack{\mathcal{C}\!\mathit{oh}}
\def\Spacesstack{\mathcal{S}\!\mathit{paces}}
\def\Quotfunctor{\mathrm{Quot}}
\def\Hilbfunctor{\mathrm{Hilb}}
\def\Curvesstack{\mathcal{C}\!\mathit{urves}}
\def\Polarizedstack{\mathcal{P}\!\mathit{olarized}}
\def\Complexesstack{\mathcal{C}\!\mathit{omplexes}}
% \Pic is the operator that assigns to X its picard group, usage \Pic(X)
% \Picardstack_{X/B} denotes the Picard stack of X over B
% \Picardfunctor_{X/B} denotes the Picard functor of X over B
\def\Pic{\mathop{\mathrm{Pic}}\nolimits}
\def\Picardstack{\mathcal{P}\!\mathit{ic}}
\def\Picardfunctor{\mathrm{Pic}}
\def\Deformationcategory{\mathcal{D}\!\mathit{ef}}
\usepackage{cleveref}
\crefname{originalreference}{original reference}{original references}
\Crefname{originalreference}{Original reference}{Original references}
\makeatletter
\newlabel{definition-set-theoretically-invariant}{{048V}{source}{}{sourcepointer-1}{}}
\newlabel{definition-set-theoretically-invariant@cref}{{[originalreference][0][]048V}{[1][0][]source}{}{}{}}
\makeatother
\crefname{originalreference}{original reference}{original references}
\Crefname{originalreference}{Original reference}{Original references}
\makeatletter
\newlabel{lemma-geometric-orbits}{{048U}{source}{}{sourcepointer-2}{}}
\newlabel{lemma-geometric-orbits@cref}{{[originalreference][0][]048U}{[1][0][]source}{}{}{}}
\makeatother
\crefname{originalsection}{original section}{original sections}
\Crefname{originalsection}{Original section}{Original sections}
\makeatletter
\newlabel{section-conventions-notation}{{048C}{source}{}{sourcepointer-3}{}}
\newlabel{section-conventions-notation@cref}{{[originalsection][0][]048C}{[1][0][]source}{}{}{}}
\makeatother
\crefname{originallemma}{original lemma}{original lemmas}
\Crefname{originallemma}{Original lemma}{Original lemmas}
\makeatletter
\newlabel{spaces-morphisms-lemma-locally-finite-type-surjective-geometric-points}{{0487}{source}{}{sourcepointer-4}{}}
\newlabel{spaces-morphisms-lemma-locally-finite-type-surjective-geometric-points@cref}{{[originallemma][0][]0487}{[1][0][]source}{}{}{}}
\makeatother
\crefname{originallemma}{original lemma}{original lemmas}
\Crefname{originallemma}{Original lemma}{Original lemmas}
\makeatletter
\newlabel{spaces-morphisms-lemma-permanence-finite-type}{{0462}{source}{}{sourcepointer-5}{}}
\newlabel{spaces-morphisms-lemma-permanence-finite-type@cref}{{[originallemma][0][]0462}{[1][0][]source}{}{}{}}
\makeatother
\title{Orbit spaces over a locally finite type pre-relation detected on k-valued points}
\author{The Stacks Project Authors}
\date{stacks-project:groupoids-quotients.tex:local-snapshot-20261012; GNU FDL 1.2 or later}
\begin{document}
\maketitle
\noindent\textit{Editorial note.} All delivered mathematical text is the author's own, copied from the pinned Stacks Project chapter file groupoids-quotients.tex; no statement or proof step was rewritten, reordered or removed.\par\medskip
\noindent\textit{Editorial note.} The counted proof is the author's complete argument: This is the pointwise reincarnation of the orbit space definition under the locally finite type hypothesis.\par\medskip
\noindent\textit{Editorial note.} The label lemma-orbit-space-locally-finite-type-over-base names the counted statement, and the proof environment immediately following it in the source is delivered verbatim as its proof.\par\medskip
\noindent\textit{Editorial note.} The counted statement is carried by the author's own lemma environment in the chapter source.\par\medskip
\noindent\textit{Allowed standard prerequisites (editorial).}\begin{itemize}
\item Chapter Definition definition-orbit-space (navigation stub): phi is an orbit space for R when it is R-invariant, separates R-orbits, and is surjective
\item Chapter Definition definition-set-theoretically-invariant (navigation stub): set-theoretically R-invariant means U(k) to X(k) equalizes s and t for every algebraically closed k over B; separating orbits means equal values force the points into the same orbit
\item Delivered card X2-alg-groupoids-quotients-geometric-orbits (Lemma lemma-geometric-orbits): under the locally finite type hypothesis, being in the same R-orbit on k-valued points agrees with equivalence under the relation generated by j(R(k))
\item Morphisms of Spaces, Lemma spaces-morphisms-lemma-permanence-finite-type: local finite type is preserved under composition
\item Morphisms of Spaces, Lemma spaces-morphisms-lemma-locally-finite-type-surjective-geometric-points: a nonempty algebraic space locally of finite type over an algebraically closed field has a k-valued point
\end{itemize}\medskip
\ifdefined\hypertarget\hypertarget{sourcepointer-1}{}\fi
\noindent\textit{Source reference.} Original reference 048V: sources/stacks-groupoids-quotients/pinned-original-source.tex, source lines 953--953. Prerequisite pointer invoked at 1429; it names a standard fact of the same chapter that the counted argument uses, and every step of the counted proof built on it is retained verbatim in the delivered body.\par\medskip
\ifdefined\hypertarget\hypertarget{sourcepointer-2}{}\fi
\noindent\textit{Source reference.} Original reference 048U: sources/stacks-groupoids-quotients/pinned-original-source.tex, source lines 892--892. Prerequisite pointer invoked at 1431; it names a standard fact of the same chapter that the counted argument uses, and every step of the counted proof built on it is retained verbatim in the delivered body.\par\medskip
\ifdefined\hypertarget\hypertarget{sourcepointer-3}{}\fi
\noindent\textit{Source reference.} Original section 048C: sources/stacks-groupoids-quotients/pinned-original-source.tex, source lines 249--249. Prerequisite pointer invoked at 1401; it names a standard fact of the same chapter that the counted argument uses, and every step of the counted proof built on it is retained verbatim in the delivered body.\par\medskip
\ifdefined\hypertarget\hypertarget{sourcepointer-4}{}\fi
\noindent\textit{Source reference.} Original lemma 0487: sources/stacks-chow/tags.txt, source lines 5534--5534. Prerequisite pointer invoked at 1420; it names a standard fact of a sibling chapter that the counted argument uses, and every step of the counted proof built on it is retained verbatim in the delivered body.\par\medskip
\ifdefined\hypertarget\hypertarget{sourcepointer-5}{}\fi
\noindent\textit{Source reference.} Original lemma 0462: sources/stacks-chow/tags.txt, source lines 5457--5457. Prerequisite pointer invoked at 1417; it names a standard fact of a sibling chapter that the counted argument uses, and every step of the counted proof built on it is retained verbatim in the delivered body.\par\medskip
\begin{lemma}
\label{lemma-orbit-space-locally-finite-type-over-base}
Let $B \to S$ as in Section \ref{section-conventions-notation}.
Let $j = (t, s) : R \to U \times_B U$ be a pre-relation.
Assume $R, U$ are locally of finite type over $B$.
Let $\phi : U \to X$ be an $R$-invariant morphism of algebraic spaces over $B$.
Then $\phi$ is an orbit space for $R$ if and only if the natural map
$$
U(k)/\big(\text{equivalence relation generated by }j(R(k))\big)
\longrightarrow
X(k)
$$
is bijective for all algebraically closed fields $k$ over $B$.
\end{lemma}

\begin{proof}
Note that since $U$, $R$ are locally of finite type over $B$ all of the
morphisms $s, t, j, \phi$ are locally of finite type, see
Morphisms of Spaces, Lemma \ref{spaces-morphisms-lemma-permanence-finite-type}.
We will also use without further mention
Morphisms of Spaces, Lemma
\ref{spaces-morphisms-lemma-locally-finite-type-surjective-geometric-points}.
Assume $\phi$ is an orbit space. Let $k$ be any algebraically closed
field over $B$. Let $\overline{x} \in X(k)$. Consider
$U \times_{\phi, X, \overline{x}} \Spec(k)$.
This is a nonempty algebraic space
which is locally of finite type over $k$. Hence it has a $k$-valued point.
This shows the displayed map of the lemma is surjective.
Suppose that $\overline{u}, \overline{u}' \in U(k)$ map to the same
element of $X(k)$. By
Definition \ref{definition-set-theoretically-invariant}
this means that $\overline{u}, \overline{u}'$ are in the same
$R$-orbit. By Lemma \ref{lemma-geometric-orbits} this means that
they are equivalent under the equivalence relation generated by
$j(R(k))$. Thus the displayed morphism is injective.

\medskip\noindent
Conversely, assume the displayed map is bijective for all algebraically
closed fields $k$ over $B$. This condition clearly implies that $\phi$
is surjective. We have already assumed that $\phi$ is $R$-invariant.
Finally, the injectivity of all the displayed maps implies that
$\phi$ separates orbits. Hence $\phi$ is an orbit space.
\end{proof}

\end{document}
